\section*{Chapter \ref{ch:b3:social-insurance} --- Social Insurance}

\begin{solution}{exo:b3:social-insurance:1}
From 0.09 weeks (Austrian men, a lower bound) to 0.32 weeks (Austrian women) of time
out of work per week of entitlement; time on benefits rises by about 0.5 weeks in
Missouri. Benefit receipt rises more because people who would have found work
after their benefits ran out claim the extra weeks, and some who are working
informally or about to start a job draw benefits meanwhile. % figdata: duration
\end{solution}

\begin{solution}{exo:b3:social-insurance:2}
$0.6 \times 2{,}500 = 1{,}500$: a replacement rate of 60~\%. At 4{,}000, $0.6 \times 4{,}000 = 2{,}400$ is capped at 1{,}800: a rate of
45~\%.
\end{solution}

\begin{solution}{exo:b3:social-insurance:3}
$0.5 + 1.5 = 2~\%$ a year. With a population shrinking by 0.5~\%, $-0.5 + 1.5 = 1~\%$.
\end{solution}

\begin{solution}{exo:b3:social-insurance:4}
No jump in the rate of finding work when benefits run out; no difference in spell
length at a threshold where entitlement changes; or durations that respond as much
for people whose benefits do not change as for those whose benefits do.
\end{solution}

\begin{solution}{exo:b3:social-insurance:5}
Gain $2 \times 0.10 = 0.20$; cost $0.5/0.95 = 0.526$. The gain is below the cost: benefits should fall.
They are optimal when $2\,\Delta c/c = 0.526$, a consumption drop of 26.3~\%.
\end{solution}

\begin{solution}{exo:b3:social-insurance:6}
$13 \times 0.16 = 2.1$ to $13 \times 0.20 = 2.6$ weeks longer.
\end{solution}

\begin{solution}{exo:b3:social-insurance:7}
\omterm{def:b3:social-insurance:annuity}{Annuity}: $0.8/1.03 = 0.777$; bond: $1/1.03 = 0.971$. The \omterm{def:b3:social-insurance:annuity}{annuity} is cheaper because it pays nothing
in the 20~\% of cases where the retiree has died.
\end{solution}

\begin{solution}{exo:b3:social-insurance:8}
The tested solution is \code{ex08_ageing.py}, in the chapter's \code{exercises/}
folder. Among the 160 countries of a million people or more with both years, the
ratio rose most in Japan (25.6 to 51.0), Puerto Rico (17.9 to 39.8), Hong Kong (15.5 to
35.8), Bosnia and Herzegovina (15.9 to 35.6) and Korea (9.9 to 29.3).
% output: ex08_ageing
\end{solution}

\begin{solution}{exo:b3:social-insurance:9}
Significantly lower blood pressure, cholesterol or blood sugar among those drawn
than among those not drawn, scaled up by the 25-point difference in coverage;
the measured differences were not significant.
\end{solution}

\begin{solution}{exo:b3:social-insurance:10}
The proof is in the text. The envelope theorem enters when $W$ is differentiated in
$b$: the worker's own utility does not change to first order through $D$, since she
chose $D$ optimally, so $dD/db$ matters only through the fund's budget. The search cost
$\psi$ enters only through the choice of $D$, which the measured elasticity $\varepsilon$ summarises.
\end{solution}

\begin{solution}{exo:b3:social-insurance:11}
At the pooled price of 0.7 per dalen paid in period 2, short-lived retirees value the
\omterm{def:b3:social-insurance:annuity}{annuity} at 0.5 and do not buy; only long-lived ones buy, and the insurer, paying 0.9
per dalen sold, loses money. The price rises to 0.9, where only the long-lived buy
and the insurer breaks even: the market has unravelled to one type.
\end{solution}

\begin{solution}{exo:b3:social-insurance:12}
The tested solution is \code{ex12_age_rd.py}, in the chapter's \code{exercises/}
folder. With 9{,}965 observations within five years of 50, the jump is 0.769 months
(s.e.\ 0.119), so 0.128 months per month of entitlement, against a true 0.150: the
95~\% interval of the jump, 0.54 to 1.00, contains the true 0.9. % output: ex12_age_rd
\end{solution}

\begin{solution}{pb:b3:social-insurance:fund}
% figdata: fund_stats, fund_curve
\textbf{1.} $4{,}645 - 4{,}366 = 279$ thousand, 6.0~\%.

\textbf{2.} $0.6 \times 31{,}000 = 18{,}600$ dalens.

\textbf{3.} $279{,}000 \times 18{,}600 = 5{,}189$ million dalens.

\textbf{4.} $t = 0.6 \times 0.06/0.94 = 3.83~\%$ of wages.

\textbf{5.} $0.24 - 0.28 \times 0.6 = 7.2~\%$; with no insurance, 24~\%.

\textbf{6.} Many have little saved and cannot borrow against future earnings while
unemployed; a long spell exhausts savings.

\textbf{7.} $(35/60)^{0.4} - 1 = -19.3~\%$.

\textbf{8.} $6.0 \times 0.807 = 4.8~\%$.

\textbf{9.} It falls: fewer people draw benefits, for less time. Each unemployed person's
choice of how long to search changes what every worker pays, a cost she does not
bear herself: a fiscal externality.

\textbf{10.} Benefits lower the cost of searching less hard or of turning offers down
(moral hazard), and they relieve people without savings of the need to take the
first job (liquidity).

\textbf{11.} A payment that does not depend on staying unemployed also lengthens
spells, so part of the response to benefits is liquidity, not moral hazard.

\textbf{12.} The liquidity part is a gain to the unemployed, who search as they would
with savings; the formula's cost should count only the moral-hazard part.

\textbf{13.} $3 \times 0.072 = 0.216$.

\textbf{14.} $0.4/0.94 = 0.426$.

\textbf{15.} Too generous: the gain is below the cost.

\textbf{16.} $3(0.24 - 0.28r) = 0.426$ gives $r = (0.24 - 0.142)/0.28 = 35.1~\%$.

\textbf{17.} $0.351 \times 0.048/0.952 = 1.77~\%$, against 3.83~\% at 60~\%. With unemployment
held at 6~\% it would be $0.351 \times 0.06/0.94 = 2.24~\%$: of the fall of 2.06 points, 0.47, almost a
quarter, comes from lower unemployment --- the fiscal externality of question 9.

\textbf{18.} $0.24 - 0.28 \times 0.351 = 14.2~\%$.

\textbf{19.} $(1 - \Delta c/c)^{-3} - 1 = 0.426$ gives a drop of 11.2~\% and $r = 45.9~\%$. The exact gain $(1 - d)^{-3} - 1$ exceeds $3d$,
increasingly so for large drops, so insurance is worth more than the approximation
says.

\textbf{20.} $(0.24 - 0.213/3)/0.28 = 60.4~\%$.

\textbf{21.} $(0.24 - 0.213)/0.28 = 9.7~\%$.

\textbf{22.} $0.9/6 = 0.15$ months per month.

\textbf{23.} Both are ratios of time to time, and 0.15 lies inside the range of 0.09 to 0.32.

\textbf{24.} Workers cannot choose their age, and those just below and just above 50
are alike in other ways; the estimate applies to job losers near 50 only, and a
duration elasticity is not the benefit-level elasticity the formula uses.

\textbf{25.} With an elasticity of 0.4, relative risk aversion 3 and the survey's
consumption drops, Brenland's optimal replacement rate is 35~\% by the approximate
formula (46~\% by the exact one), below its 60~\%.
\end{solution}
